\chapter{Code reproduction guide}\label{ch:appendix-user-guide}

This appendix explains how to reproduce every table and figure in this report from the source release. All commands assume the repository root as the working directory.

\section*{Environment}

The code targets CPU and is known to run on a laptop. Python $\ge 3.10$ is required.

\begin{verbatim}
python3 -m venv .venv
source .venv/bin/activate
pip install --upgrade pip
pip install -r requirements.txt
\end{verbatim}

Alternatively, under Anaconda:
\begin{verbatim}
conda create -n srfca python=3.11
conda activate srfca
pip install -r requirements.txt
\end{verbatim}

The pinned requirements cover NumPy, SciPy, scikit-learn, PyTorch (CPU), torchvision, networkx, pandas, matplotlib, seaborn, jupyter and tqdm.

\section*{Smoke test}

Before running any MNIST experiment, confirm that SR-FCA's internals are wired correctly by running the synthetic 2-Gaussian sanity check:

\begin{verbatim}
python -m src.srfca --smoke-test --seed 0
\end{verbatim}

Expected output includes \texttt{misclustering\_error = 0.0000} and \texttt{recovered 2 clusters}. This finishes in under 30 seconds on CPU.

\section*{Reproducing the report's figures and tables}

Every table and figure in Chapters~\ref{ch:experiments} and~\ref{ch:part-c} is regenerated by running four Jupyter notebooks in order:

\begin{verbatim}
jupyter nbconvert --to notebook --execute --inplace \
    notebooks/01_main_experiments.ipynb
jupyter nbconvert --to notebook --execute --inplace \
    notebooks/02_ablation_study.ipynb
jupyter nbconvert --to notebook --execute --inplace \
    notebooks/03_new_dataset.ipynb
jupyter nbconvert --to notebook --execute --inplace \
    notebooks/04_application.ipynb
\end{verbatim}

Outputs land in:
\begin{itemize}
    \item \texttt{results/*.csv} -- raw numeric results, one row per (dataset, seed) cell of each notebook's grid.
    \item \texttt{results/figures/*.png} -- every figure included in the report.
    \item \texttt{results/cache/*.pkl} -- per-cell cache for resumability; do not delete if you intend to re-run.
\end{itemize}

The notebooks are resumable: each (dataset, seed) cell is cached after the first successful run, so interrupting and re-running picks up where it stopped. This is important on CPU where the four notebooks together take ${\approx}13$ hours of wall clock.

\section*{Mapping from results to report objects}

\begin{itemize}
    \item \textbf{Tables~\ref{tab:main-acc}, \ref{tab:main-mis}, Figure~\ref{fig:main-bar}} $\leftarrow$ \texttt{results/main.csv}, \texttt{main\_bar.png}.
    \item \textbf{Figure~\ref{fig:pairwise-hist}} $\leftarrow$ \texttt{pairwise\_hist\_\{rotated,inverted\}.png}.
    \item \textbf{Table~\ref{tab:ablation-A1}, Figure~\ref{fig:ablation-A1}} $\leftarrow$ \texttt{results/ablation\_A1.csv}, \texttt{ablation\_A1.png}.
    \item \textbf{Table~\ref{tab:fashion}, Figure~\ref{fig:fashion-bar}} $\leftarrow$ \texttt{results/fashion\_mnist.csv}, \texttt{fashion\_mnist\_bar.png}, \texttt{pairwise\_hist\_fashion\_mnist.png}.
    \item \textbf{Table~\ref{tab:part-c}, Figure~\ref{fig:part-c}} $\leftarrow$ \texttt{results/part\_c.csv}, \texttt{part\_c\_bar.png}, \texttt{pairwise\_hist\_part\_c.png}.
\end{itemize}

\section*{Source-code layout}

\begin{itemize}
    \item \texttt{src/datasets.py} -- MNIST and Fashion-MNIST federated splitters (rotated/inverted); synthetic Gaussian mixtures for the smoke test.
    \item \texttt{src/models.py} -- 2-layer MLP and a linear-regression head, both exposing flat-weight read/write for aggregation.
    \item \texttt{src/federated.py} -- local SGD, FedAvg, coordinate-wise trimmed mean, and the \texttt{train\_cluster\_model} routine used by SR-FCA's REFINE.
    \item \texttt{src/srfca.py} -- all four subroutines (ONE\_SHOT, TrimmedMeanGD, RECLUSTER, MERGE), the outer loop \texttt{sr\_fca}, and the synthetic smoke test.
    \item \texttt{src/baselines.py} -- Local, Global (FedAvg), IFCA, Local-KMeans.
    \item \texttt{src/metrics.py} -- Hungarian-matching test accuracy and misclustering error (the two metrics reported in the main text); NMI and ARI helpers are also defined for completeness but no longer consumed by the notebooks after the metric set was aligned to the paper.
    \item \texttt{src/utils.py} -- seeding, device selection, and path helpers.
\end{itemize}

\section*{Reproducibility notes}

\begin{itemize}
    \item Global seed $42$ is set via \texttt{src.utils.set\_seed}; this seeds Python's \texttt{random}, NumPy, and PyTorch CPU generators.
    \item MNIST and Fashion-MNIST are downloaded via \texttt{torchvision.datasets} into \texttt{./data/} on first run; subsequent runs use the cached files.
    \item Every random choice in the experiments (client shuffling, cluster assignment, batch order) is a function of the seed only; re-running the same notebook produces bit-identical CSVs on the same machine.
\end{itemize}
